\section{Introduction}

The motivating prime-shell collision is
\begin{equation}
 7p+3r=16Q,\qquad Q<p<r<2Q,                         \label{eq:resonance}
\end{equation}
with the primitive condition $(Q,21)=1$.  Earlier exact work shows that these
collisions form a matching and that a proportional energy saving cannot exceed the
diagonal mass carried by its incident vertices.  In the literal aligned row model,
strict saving by $\delta$ therefore requires
\begin{equation}
 \frac{E_{3,7}(Q)}{|\PP_Q|}\ge \frac{\delta}{8}.     \label{eq:toll}
\end{equation}
For the program endpoint $\delta=1/400$, the toll is $1/3200$.

The purpose of this paper is to decide whether the left side of
\eqref{eq:toll} can remain bounded below.  The answer is no.  The relevant arithmetic
is an ordinary dimension-two upper-bound sieve, but its application requires keeping
the moving determinant $16Q$ and its exceptional local factors.  We use the classical
Selberg sieve in the forms developed in \cite{HalberstamRichert1974,IwaniecKowalski2004};
the Goldbach and twin-prime derivations in \cite{GreenAdditive,Rudnick2015} provide
close two-form templates.

Our second result isolates the real scope of the obstruction: any \emph{fixed finite}
set of nondegenerate linear resonance equations has zero edge density in the prime
shell.  A family whose depth grows with $Q$, or an arithmetic concentration of actual
source mass on a sparse support, is not ruled out.
